\clearpage
\section*{Problem 8}
\addcontentsline{toc}{section}{Problem 8}
\begingroup
\hypersetup{colorlinks=true,
            linkcolor=blue!50!black,
            citecolor=blue!50!black,
            urlcolor=blue!50!black}
\newcommand{\qeightR}{\mathbb{R}}
\newcommand{\qeightZ}{\mathbb{Z}}
\newcommand{\qeightSm}{\mathscr S}
\newcommand{\qeightLdual}{L^{\vee}}
\newcommand{\qeightdd}{\mathrm d}
\newcommand{\qeightLag}{\Lambda}
\newcommand{\qeightHess}{\operatorname{Hess}}
\newtheorem{qeighttheorem}{Theorem}
\newtheorem{qeightproposition}{Proposition}
\newtheorem{qeightlemma}{Lemma}
\newtheorem{qeightcorollary}{Corollary}
\theoremstyle{definition}
\newtheorem{qeightdefinition}{Definition}
\newtheorem{qeightremark}{Remark}
\title{Smoothing Polyhedral Lagrangian Surfaces}
\author{}
\date{}
\maketitle


\section*{Statement and Answer}

Equip $\qeightR^4$ with coordinates $(q_1,q_2,p_1,p_2)$ and the standard symplectic
form $\omega=\qeightdd p_1\wedge \qeightdd q_1+\qeightdd p_2\wedge \qeightdd q_2$. A \emph{polyhedral
Lagrangian surface} $K\subset\qeightR^4$ is a finite polyhedral complex, all of whose
faces are Lagrangian, that is a topological submanifold of $\qeightR^4$.

\medskip
\noindent\textbf{Theorem.} \emph{If $K$ is a polyhedral Lagrangian surface in
which exactly $4$ faces meet at every vertex, then $K$ admits a Lagrangian
smoothing: a Hamiltonian isotopy $K_t$ of smooth Lagrangian submanifolds
parametrized by $t\in(0,1]$, extending to a topological isotopy on $[0,1]$ with
$K_0=K$.}

\medskip
\noindent The answer is \textbf{YES}.

\medskip
The proof has a local part (a normal form near each vertex and edge, from which
a smooth structure on $K$ is read off) and a global part (a fibration of
Lagrangian planes transverse to $K$, with respect to which all sufficiently
$C^1$-small ``smoothing functions'' yield genuine graphical Lagrangians). The
$4$-faces hypothesis enters in the local normal form (Lemma~\ref{lem:linear}),
where four Lagrangian quadrants force a standard symplectic basis; it is not a
Maslov/parity condition.

\section*{1. Local model at a vertex}

\begin{qeightlemma}\label{lem:linear}
Let $\qeightR^2\to\qeightR^4$ be an embedding that is linear on each of the four quadrants,
with Lagrangian image $\Sigma$, and assume $\Sigma$ is not contained in a single
plane. Then there is a linear symplectic transformation of $\qeightR^4$ carrying
$\Sigma$ to the product of the union of the positive coordinate axes in
$\qeightR^2_{p_1q_1}$ with that in $\qeightR^2_{p_2q_2}$.
\end{qeightlemma}

\begin{proof}
Let $(v_1,v_2,u_1,u_2)$ be the tangent vectors at the origin to the four edges of
$\Sigma$, ordered so that cyclically adjacent vectors span the faces. The
pairings $\omega(v_i,u_i)$ cannot vanish: if some $\omega(v_i,u_i)=0$, then
$\omega$ would vanish identically on a $3$-dimensional subspace of $\qeightR^4$, which
is impossible for a symplectic form. After swapping $v_i\leftrightarrow u_i$ if
needed we may take both pairings positive, and after rescaling we may take
$\omega(v_1,u_1)=\omega(v_2,u_2)=1$ with the cross pairings vanishing (the faces
are Lagrangian). Then $(v_1,v_2,u_1,u_2)$ is a standard symplectic basis, and the
linear map $\partial_{p_i}\mapsto v_i$, $\partial_{q_i}\mapsto u_i$ is the
required symplectomorphism, carrying $\Sigma$ to the stated product of
positive-axis unions.
\end{proof}

In the plane $\qeightR^2_{pq}$ the symplectic pairing projects the union of the two
positive axes homeomorphically onto the dual of the antidiagonal line $p=q$.
Taking the product and applying Lemma~\ref{lem:linear}:

\begin{qeightcorollary}\label{cor:Ldual}
There is a linear Lagrangian plane $L\subset\qeightR^4$ such that the symplectic
pairing $\qeightR^4\to \qeightLdual$ restricts to a homeomorphism $\Sigma\to \qeightLdual$.
\end{qeightcorollary}

This equips $\Sigma$ with a smooth structure (pulled back from $\qeightLdual$), which
we fix henceforth. Given $L$, call a Lagrangian $\qeightLag\subset\qeightR^4$
\emph{graphical} if the pairing restricts to a diffeomorphism $\qeightLag\cong\qeightLdual$.

\section*{2. Smoothing functions}

Choose a Lagrangian splitting $\qeightR^4\cong L\oplus\qeightLdual$. Each of the four
quadrant-faces of $\Sigma$ is graphical over $\qeightLdual$, hence the graph of the
differential of a quadratic form; collecting these gives quadratic forms
$\{q_{ij}\}_{i,j\in\pm}$ on $\qeightLdual$, agreeing to first order along the edges.
Via $\Sigma\cong\qeightLdual$ (Corollary~\ref{cor:Ldual}) write $q_\Sigma$ for the
resulting $C^1$ function on $\Sigma$ that is $q_{ij}$ on each face.

\begin{qeightdefinition}\label{def:smoothing}
The space $\qeightSm(\Sigma)$ of \emph{smoothing functions} is the set of $C^1$
functions $f:\Sigma\to\qeightR$ such that $f+q_\Sigma$ is $C^\infty$.
\end{qeightdefinition}

$\qeightSm(\Sigma)$ is invariant under adding smooth functions, and depends only on
$L$, not on the splitting (a different splitting changes $q_{ij}$ by a global
quadratic form, i.e.\ a smooth function). To a smoothing function $f$ we
associate the Lagrangian $\qeightLag_{\qeightdd f}\subset\qeightR^4$ that is, piecewise, the graph
of $\qeightdd f$ over each face — equivalently the graph of $\qeightdd(f+q_\Sigma)$ in the
splitting.

\begin{qeightlemma}\label{lem:bijection}
$f\mapsto \qeightLag_{\qeightdd f}$ is a bijection between graphical Lagrangians and
smoothing functions on $\Sigma$ modulo additive constants.
\end{qeightlemma}

\begin{proof}
In the splitting, $\qeightLag_{\qeightdd f}$ is the graph of $\qeightdd(f+q_\Sigma)$ as a function
on $\qeightLdual$; graphical Lagrangians over $\qeightLdual$ are exactly graphs of
differentials of smooth functions, and $f+q_\Sigma$ is smooth precisely when
$f\in\qeightSm(\Sigma)$. The construction does not use the splitting, so the bijection
depends only on $L$.
\end{proof}

\noindent The same statements hold verbatim for an edge: a pair of Lagrangian
half-planes meeting along a line, with the analogue of
Corollary~\ref{cor:Ldual} provided by the next lemma.

\begin{qeightlemma}\label{lem:edge}
If $\Sigma$ is a pair of linear Lagrangian half-planes meeting along a line
$\ell$, the space of Lagrangian planes $L$ for which the pairing
$\Sigma\to\qeightLdual$ is a homeomorphism is contractible.
\end{qeightlemma}

\begin{proof}
Up to affine linear symplectomorphism, $\Sigma$ is the product of the real axis
in one $\qeightR^2$-factor with the union of the positive axes in another. The pairing
$\Sigma\to\qeightLdual$ is a homeomorphism iff $L$ is transverse to both half-planes,
i.e.\ the symplectic reduction of $L$ along $\ell$ is a line in $\ell^\perp/\ell\cong\qeightR^2$
meeting the interior of the positive quadrant — a contractible condition. The
space of Lagrangian lifts of such a line is itself contractible (two lifts differ
by the graph of a map $\ell'\to\ell$), giving the claim.
\end{proof}

\section*{3. Existence of small smoothing functions}

The crux is that $\qeightSm(\Sigma)$ is \emph{not} invariant under rescaling, so the
existence of small smoothing functions is not automatic.

\begin{qeightlemma}\label{lem:small}
$\qeightSm(\Sigma)$ contains functions of arbitrarily small $C^1$-norm.
\end{qeightlemma}

\begin{proof}
Fix a partition of unity $\sum_\sigma \chi_\sigma=1$ on $\Sigma$ indexed by the
strata, with $\chi_\sigma$ supported near $\sigma$ and $\equiv1$ on $\sigma$ away
from its boundary, of bounded $C^k$ norms. Let $\chi_\sigma^\varepsilon$ be the
composition of $\chi_\sigma$ with dilation by $1/\varepsilon$; these are
uniformly bounded with $\|\chi_\sigma^\varepsilon\|_{C^1}=O(1/\varepsilon)$.

Choose a Lagrangian plane $\Lambda_\sigma$ containing each stratum $\sigma$ and
transverse to $L$, with associated smoothing function $f_\sigma$. The first-order
agreement of the faces along shared edges means $f_\sigma$ and $f_{\sigma'}$ agree
to first order along $\sigma\cap\sigma'$. Set $f^\varepsilon:=\sum_\sigma
\chi_\sigma^\varepsilon f_\sigma$. The partition-of-unity structure makes
$f^\varepsilon$ a smoothing function, and although $\|\nabla\chi_\sigma^\varepsilon\|$
grows like $1/\varepsilon$, it is supported where $f_\sigma-f_{\sigma'}=O(\varepsilon^2)$;
by the product rule the $C^1$-norm of $f^\varepsilon$ is therefore $O(\varepsilon)$.
\end{proof}

\section*{4. Global construction: a conormal fibration}

To globalize, we replace the single plane $L$ by a smoothly varying family.

\begin{qeightdefinition}\label{def:fibration}
A \emph{conormal fibration dual to $K$} is a smoothly varying family $L_z$ of
affine Lagrangian planes, $z\in K$, such that: (i) near each vertex the $L_z$ are
translates of a single plane as in Corollary~\ref{cor:Ldual}; (ii) $L_z$ varies
smoothly along edges and along faces; and (iii) in a collar of each edge the
planes in the normal direction are translates of those along the edge, with
\emph{matched} inward normal fields on the two adjacent faces (opposite vectors
under the identification $T_z\sigma\cong\qeightLdual_z\cong T_z\sigma'$).
\end{qeightdefinition}

\begin{qeightlemma}\label{lem:fibexists}
$K$ admits a conormal fibration which near each vertex agrees with the local
choice of Corollary~\ref{cor:Ldual}.
\end{qeightlemma}

\begin{proof}
Lemma~\ref{lem:edge} extends the vertex choices over the edges (the space of
admissible $L_z$ along an edge is contractible). Choosing a normal vector field
to one face along each edge determines matched normals on the other; extending to
the interior of the faces is unobstructed because the space of Lagrangian planes
transverse to a given one is contractible.
\end{proof}

The fibration determines an open neighbourhood $\nu K$ of $K$ in which the fibres
$L_z$ are disjoint, hence a projection $\pi\colon\nu K\to K$ (each point $w\in\nu K$
lies in a unique fibre $L_z$, and $\pi(w):=z$). A Lagrangian contained in
$\nu K$ is \emph{graphical} (with respect to $L_z$) if its projection $\pi$ to $K$
is a diffeomorphism onto $K$. The local correspondence of
Lemma~\ref{lem:bijection} globalizes:

\begin{qeightlemma}\label{lem:graph-global}
Every graphical Lagrangian with respect to $\{L_z\}$ is the graph of a smoothing
function on $K$; and every smoothing function of sufficiently small differential
defines a graphical Lagrangian.
\end{qeightlemma}

\begin{proof}
The correspondence is local on $K$. At a point $z$, a plane $\qeightLdual_z$ transverse
to $L_z$ stays transverse to nearby fibres, so a neighbourhood of $z$ in $\nu K$
is modelled on the conormal bundle of $\qeightLdual_z$ (Weinstein's tubular
neighbourhood theorem), and Lagrangians there are graphs of closed $1$-forms,
i.e.\ differentials of smoothing functions.
\end{proof}

\begin{qeightlemma}\label{lem:small-global}
There exist smoothing functions for $K$ of arbitrarily small $C^1$-norm.
\end{qeightlemma}

\begin{proof}
Take a partition of unity $\sum_\alpha\rho_\alpha=1$ on $K$ indexed by strata,
with $\rho_\alpha$ supported in the open star of $\alpha$. Lemma~\ref{lem:small}
gives smoothing functions $f_\alpha$ of arbitrarily small $C^1$-norm on each open
star; then $\sum_\alpha\rho_\alpha f_\alpha$ is a global smoothing function of
small $C^1$-norm.
\end{proof}

\section*{5. From smoothing functions to sections of the fibration}

We now record the precise relationship between graphical Lagrangians and
\emph{embeddings of $K$ into $\qeightR^4$}, which is what is required to produce a
topological isotopy. The point is that the conormal fibration provides a fixed
parametrizing base $K$ for every graphical Lagrangian, and the graph map is
manifestly injective and continuous up to the singular limit.

\paragraph{Sections of the fibration.} Fix the conormal fibration $\{L_z\}_{z\in K}$
and its projection $\pi\colon\nu K\to K$ from
Definition~\ref{def:fibration}. A \emph{section of $\pi$} is a continuous map
$s\colon K\to\nu K$ with $\pi\circ s=\mathrm{id}_K$; thus $s(z)\in L_z$ for every
$z$. The zero section $s\equiv z$ (i.e.\ $s(z)=z\in K\subset\nu K$, using
$z\in L_z$) parametrizes $K$ itself. We measure the size of a section by
\[
\|s\|_{C^0}:=\sup_{z\in K}\,d_{\qeightR^4}\bigl(s(z),\,z\bigr),
\]
the largest Euclidean displacement of $s$ from the zero section.

\begin{qeightlemma}\label{lem:section-embed}
Let $s\colon K\to\nu K$ be a continuous section of $\pi$. Then $s$ is a
topological embedding of $K$ into $\qeightR^4$, with image $s(K)$. If $s$ is the
zero section then $s(K)=K$; if $s$ corresponds (in the sense of
Lemma~\ref{lem:graph-global}) to a smoothing function $f$ of sufficiently small
differential, then $s$ is smooth on each face and $s(K)=\qeightLag_{\qeightdd f}$.
\end{qeightlemma}

\begin{proof}
\emph{Continuity.} $s$ is continuous by assumption, as a map into $\nu K\subset\qeightR^4$.

\emph{Injectivity, uniform in the section.} Suppose $s(z)=s(z')$. Applying $\pi$
and using $\pi\circ s=\mathrm{id}_K$ gives $z=\pi(s(z))=\pi(s(z'))=z'$. Hence
$s$ is injective. This argument uses \emph{only} the identity $\pi\circ s=\mathrm{id}_K$;
it is therefore valid for \emph{every} section, with no smallness hypothesis, and
in particular holds at the singular base $K$ (the zero section) exactly as for any
smooth graphical Lagrangian.

\emph{Embedding.} $K$ is compact (a finite polyhedral complex) and $\qeightR^4$ is
Hausdorff, so the continuous injection $s$ is a homeomorphism onto its image
$s(K)$; thus $s$ is a topological embedding.

The final statements are the definitions: the zero section has image $K$ by
construction, and the section attached to a smoothing function $f$ has image the
graphical Lagrangian $\qeightLag_{\qeightdd f}$ by Lemma~\ref{lem:graph-global}, smooth on each
face because $f$ is.
\end{proof}

\begin{qeightlemma}\label{lem:section-small}
For each $\delta>0$ there is $\eta>0$ with the following property: if $f$ is a
smoothing function on $K$ with $\|\qeightdd f\|_{C^0}<\eta$ (small enough that
Lemma~\ref{lem:graph-global} applies) and $s_f\colon K\to\nu K$ is the
corresponding section, then $\|s_f\|_{C^0}<\delta$. In particular
$s_f\to s_0$ (the zero section) uniformly as $\|\qeightdd f\|_{C^0}\to0$.
\end{qeightlemma}

\begin{proof}
Work in a finite atlas of the compact base $K$ adapted to $\{L_z\}$: cover $K$ by
finitely many of the Weinstein charts of Lemma~\ref{lem:graph-global}, in each of
which $\nu K$ is modelled on the conormal bundle of a transverse plane
$\qeightLdual_z$, the fibre over a point being the corresponding cotangent fibre, and
the section $s_f$ is given by $z\mapsto\qeightdd(f+q_K)(z)$ in fibre coordinates,
where $q_K$ is the fixed $C^1$ reference function of Section~2 (which is itself
piecewise quadratic, hence has $\qeightdd q_K$ bounded on the compact $K$). On each
chart the chart map and its differential are uniformly bounded (the atlas is
finite), so there is a constant $C\ge1$, independent of $f$, with
\[
d_{\qeightR^4}\bigl(s_f(z),z\bigr)\ \le\ C\,\bigl|\qeightdd(f+q_K)(z)-\qeightdd q_K(z)\bigr|
\ =\ C\,\bigl|\qeightdd f(z)\bigr|\qquad(z\in K),
\]
the middle equality because the zero section $s_0$ corresponds to $f\equiv0$
(displacement zero). Taking the supremum over $z$ gives
$\|s_f\|_{C^0}\le C\,\|\qeightdd f\|_{C^0}$. Choosing $\eta:=\delta/C$ (and
$\eta$ no larger than the threshold of Lemma~\ref{lem:graph-global}) yields the
claim.
\end{proof}

\section*{6. Assembly of the smoothing}

The construction of Lemma~\ref{lem:small} (globalized in Lemma~\ref{lem:small-global})
produces not a single small smoothing function but a one-parameter family
$f^\varepsilon$, $\varepsilon\in(0,1]$, each member of $\qeightSm(K)$, with
$\|\qeightdd f^\varepsilon\|_{C^0}=O(\varepsilon)$. We use this family as the smoothing
path; the key extra observation is that it depends continuously (indeed smoothly)
on $\varepsilon$ and stays inside $\qeightSm(K)$ throughout, which is what lets the
section $C^0$-norm tend to $0$ while every member remains a genuine smooth
graphical Lagrangian. We never scale a smoothing function (which would leave
$\qeightSm(K)$, since $\qeightSm(K)$ is the affine coset $g+C^\infty(K)$ and $0\notin\qeightSm(K)$).

\begin{qeightlemma}\label{lem:cont-family}
There is a family $\{f^\varepsilon\}_{\varepsilon\in(0,1]}\subset\qeightSm(K)$ such
that:
\begin{enumerate}[label=\textup{(\alph*)}]
\item $\varepsilon\mapsto f^\varepsilon$ is a smooth map $(0,1]\to C^1(K)$, i.e.\
$(\varepsilon,z)\mapsto f^\varepsilon(z)$ and its $z$-derivatives are jointly
smooth on $(0,1]\times K$;
\item $\|\qeightdd f^\varepsilon\|_{C^0}\le M\varepsilon$ for a constant $M$ independent
of $\varepsilon$; and
\item for $\varepsilon$ small enough (say $\varepsilon\le\varepsilon_0$),
$\|\qeightdd f^\varepsilon\|_{C^0}<\eta_0$, the threshold of
Lemma~\ref{lem:graph-global}.
\end{enumerate}
\end{qeightlemma}

\begin{proof}
Take the global partition of unity $\sum_\alpha\rho_\alpha=1$ of
Lemma~\ref{lem:small-global} and, on each open star, the local family
$f_\alpha^\varepsilon=\sum_\sigma\chi_\sigma^\varepsilon f_\sigma$ of
Lemma~\ref{lem:small}, where $\chi_\sigma^\varepsilon(x)=\chi_\sigma(x/\varepsilon)$.
Set $f^\varepsilon:=\sum_\alpha\rho_\alpha f_\alpha^\varepsilon$. Each
$f^\varepsilon$ is a smoothing function (a finite sum of smoothing functions
weighted by smooth $\rho_\alpha$, and $\qeightSm(K)$ is closed under such
combinations), giving $f^\varepsilon\in\qeightSm(K)$. The estimate of
Lemma~\ref{lem:small}, summed against the fixed $\rho_\alpha$, gives
$\|\qeightdd f^\varepsilon\|_{C^0}\le M\varepsilon$, which is (b); (c) follows for
$\varepsilon_0:=\min(1,\eta_0/M)$. For (a), the map $(\varepsilon,x)\mapsto
\chi_\sigma(x/\varepsilon)$ is jointly smooth on $(0,1]\times K$ (composition of
the smooth $\chi_\sigma$ with the smooth $(\varepsilon,x)\mapsto x/\varepsilon$,
$\varepsilon>0$), so each $f_\alpha^\varepsilon$, and hence $f^\varepsilon$, is
jointly smooth in $(\varepsilon,z)$ on $(0,1]\times K$; this is (a).
\end{proof}

\begin{proof}[Proof of the Theorem]
By Lemma~\ref{lem:fibexists} fix a conormal fibration $\{L_z\}_{z\in K}$ with
disjoint fibres on a neighbourhood $\nu K$ and projection $\pi\colon\nu K\to K$.
Let $\eta_0>0$ be the threshold of Lemma~\ref{lem:graph-global} and take the
family $\{f^\varepsilon\}_{\varepsilon\in(0,\varepsilon_0]}$ of
Lemma~\ref{lem:cont-family}, so that each $\qeightLag_{\qeightdd f^\varepsilon}\subset\nu K$
is a smooth embedded graphical Lagrangian. Reparametrize $\varepsilon$ to the
interval $(0,1]$ by a fixed orientation-preserving diffeomorphism
$t\mapsto\varepsilon(t)$ of $(0,1]$ onto $(0,\varepsilon_0]$ with
$\varepsilon(t)\to0$ as $t\to0^+$ and $\varepsilon(t)/t$ bounded (e.g.\
$\varepsilon(t)=\varepsilon_0 t$); write $f_t:=f^{\varepsilon(t)}$ and let
$s_t:=s_{f_t}\colon K\to\nu K$ be the corresponding section
(Lemma~\ref{lem:section-embed}).

\medskip
\noindent\textbf{Step 1: the embeddings $\phi_t$ and their endpoint.}
For $t\in(0,1]$ set $\phi_t:=s_t\colon K\hookrightarrow\qeightR^4$, and let
$\phi_0:=s_0$ be the zero section of $\pi$, i.e.\ $\phi_0=\mathrm{id}_K$ with image
$K$. By Lemma~\ref{lem:section-embed} every $\phi_t$ ($t\in[0,1]$) is a
topological embedding; injectivity holds \emph{for every $t$, including $t=0$}, by
the one-line argument $\pi\circ\phi_t=\mathrm{id}_K$, which needs no smallness and
is valid at the singular base. Thus $\{\phi_t\}_{t\in[0,1]}$ is a family of
topological embeddings of $K$ into $\qeightR^4$ with $\phi_0$ the inclusion of $K$.

\medskip
\noindent\textbf{Step 2: continuity of $t\mapsto\phi_t$ on $(0,1]$.}
The space of topological embeddings of the compact $K$ into $\qeightR^4$ carries the
$C^0$ (uniform) topology, with distance
$\rho(\phi,\psi)=\sup_{z\in K}d_{\qeightR^4}(\phi(z),\psi(z))$. In each Weinstein chart
of Lemma~\ref{lem:graph-global} the section displacement of $s_{f}$ over $z$ is
$\qeightdd f(z)$ in fibre coordinates, depending continuously on $f$ in $C^1$; together
with the finite-atlas Lipschitz constant $C$ of Lemma~\ref{lem:section-small} this
gives, for $t,t'\in(0,1]$,
\begin{equation}\label{eq:lip}
\rho(\phi_t,\phi_{t'})
=\sup_{z\in K}d_{\qeightR^4}\!\bigl(s_{f_t}(z),s_{f_{t'}}(z)\bigr)
\ \le\ C\,\|\qeightdd f_t-\qeightdd f_{t'}\|_{C^0}.
\end{equation}
By Lemma~\ref{lem:cont-family}(a) and the smoothness of $t\mapsto\varepsilon(t)$,
the map $t\mapsto\qeightdd f_t=\qeightdd f^{\varepsilon(t)}$ is continuous (in fact smooth)
into $C^0(K)$ on $(0,1]$; hence \eqref{eq:lip} makes $t\mapsto\phi_t$ continuous on
$(0,1]$.

\medskip
\noindent\textbf{Step 3: continuity at the endpoint $t=0$.}
For $t\in(0,1]$, since $\phi_0=\mathrm{id}_K$ is the zero section,
Lemma~\ref{lem:section-small} (with the constant $C$) gives
\[
\rho(\phi_t,\phi_0)\ =\ \|s_{f_t}\|_{C^0}\ \le\ C\,\|\qeightdd f_t\|_{C^0}
\ \le\ C\,M\,\varepsilon(t)\ \xrightarrow[t\to0^+]{}\ 0 ,
\]
using Lemma~\ref{lem:cont-family}(b) and $\varepsilon(t)\to0$. Therefore
$\phi_t\to\phi_0$ uniformly as $t\to0^+$. Combined with Step~2, the map
$t\mapsto\phi_t$ is continuous on the \emph{closed} interval $[0,1]$, taking
values in topological embeddings of $K$, with $\phi_0=\mathrm{id}_K$. By
definition this is a \emph{topological isotopy} on $[0,1]$ with endpoint the
inclusion of $K$. This settles the endpoint claim that was previously only
asserted: the family is continuous and injective \emph{up to and including}
$t=0$, and its limit is the homeomorphism onto the singular set $K$, not merely a
Hausdorff limit of images.

\medskip
\noindent\textbf{Step 4: the Hamiltonian property on $(0,1]$.}
Put $K_t:=\phi_t(K)=\qeightLag_{\qeightdd f_t}$. For $t\in(0,1]$, $f_t+q_K$ is smooth
(because $f_t\in\qeightSm(K)$), and by Lemma~\ref{lem:cont-family}(a) the path
$t\mapsto f_t+q_K=f^{\varepsilon(t)}+q_K$ is a smooth path of smooth generating
functions on $(0,1]$, with smooth time-derivative
$\dot u_t:=\partial_t\bigl(f_t+q_K\bigr)=\partial_t f_t$ (a smooth function on
$K$, since $q_K$ is $t$-independent). On each Weinstein chart $K_t$ is the graph
of $\qeightdd(f_t+q_K)$; the smooth path of differentials of generating functions is a
path of exact Lagrangians, and by the standard correspondence between smooth paths
of generating functions and Hamiltonian isotopies
\cite[Ch.~9]{McDuffSalamon}, the function $\dot u_t$, transported to a compact
neighbourhood of $\nu K$ in $\qeightR^4$ via the (finitely many) Weinstein charts and a
cutoff, defines a smooth compactly supported Hamiltonian $H_t$ on $\qeightR^4$ whose
time-$t$ flow carries $K_{t_0}$ to $K_t$ for $t_0,t\in(0,1]$. Hence
$\{K_t\}_{t\in(0,1]}$ is a Hamiltonian isotopy of smooth Lagrangian submanifolds.

\medskip
\noindent\textbf{Conclusion.} The family $\{K_t=\phi_t(K)\}_{t\in[0,1]}$ is a
Hamiltonian isotopy of smooth Lagrangian submanifolds for $t\in(0,1]$ (Step~4),
extending to a topological isotopy $\phi_t$ on the closed interval $[0,1]$
(Steps~1–3), continuous and injective \emph{including} $t=0$, with endpoint
$\phi_0=\mathrm{id}_K$, i.e.\ $K_0=K$. This is precisely a Lagrangian smoothing of
$K$, completing the proof.
\end{proof}

\begin{qeightremark}
The decisive structural points repaired here are two. First, the conormal
fibration supplies a \emph{single fixed base} $K$ with projection $\pi$, so each
smoothing — and the singular limit — is a \emph{section} of $\pi$, hence an
embedding parametrized by the \emph{same} $K$; injectivity of every $\phi_t$ is
then the tautology $\pi\circ\phi_t=\mathrm{id}_K$, valid uniformly in $t$ and in
particular at $t=0$. Second, the path is the explicit smoothing family
$f^\varepsilon\in\qeightSm(K)$ of Lemma~\ref{lem:cont-family}, which stays inside the
admissible class $\qeightSm(K)$ for all $\varepsilon>0$ (we do \emph{not} scale a
smoothing function, since $\qeightSm(K)$ is not closed under scaling), depends smoothly
on $\varepsilon$, and has section $C^0$-norm $O(\varepsilon)\to0$. Continuity of
$t\mapsto\phi_t$ on $[0,1]$ then follows from \eqref{eq:lip} and the endpoint
estimate of Step~3. This is strictly stronger than the Hausdorff convergence of
\emph{images} used previously: Hausdorff convergence of sets does not by itself
yield a continuous family of embeddings nor control injectivity at the limit,
which is exactly the gap addressed.
\end{qeightremark}


\section*{7. Why the Maslov / Legendrian-filling obstruction does not apply}

It is natural to ask whether the answer is forced to be \textbf{NO} by an
obstruction of Maslov-class or Legendrian-DGA type (Chekanov--Eliashberg,
Ekholm--Honda--K\'alm\'an) attached to the link
$\Lambda_v:=K\cap S^3_\varepsilon(v)$ of a $4$-valent vertex. We explain
precisely why no such obstruction is present, and indeed why the framework that
produces it does not govern the present problem.

\paragraph{The vertex link.} In the normal form of Lemma~\ref{lem:linear} the
local model is $\Sigma=C_1\times C_2$, where $C_i\subset\qeightR^2_{p_iq_i}$ is the
corner (union of the two positive coordinate axes). The link of a corner in a
small circle is two points, so $\Lambda_v=\Sigma\cap S^3_\varepsilon(v)$ is the
join $\{2\text{ pts}\}*\{2\text{ pts}\}$, i.e.\ a $4$-cycle: topologically an
unknotted circle in $S^3_\varepsilon$, traversing the four faces in cyclic order.

\paragraph{The obstruction framework requires a compact cap; the smoothing builds
none.} The Maslov-class / DGA obstructions of the cited theory obstruct the
existence of a \emph{compact exact Lagrangian filling} of a Legendrian
$\Lambda\subset(S^3,\xi_{\mathrm{std}})$: a compact Lagrangian surface in the
symplectization (or in a Liouville filling of the contact boundary) whose only
boundary is $\Lambda$, cylindrical near the contact end. The classical Maslov-class
constraint is that such a filling $F$ must have vanishing Maslov class
$\mu_F\in H^1(F;\qeightZ)$, and the rotation/Thurston--Bennequin and DGA-augmentation
data of $\Lambda$ further constrain or obstruct $F$.

The Lagrangian smoothing constructed here is \emph{not} a filling of $\Lambda_v$
in this sense and creates no compact cap. By Lemma~\ref{lem:bijection} the
smoothed surface is, on a neighbourhood of $v$, the graph
$\qeightLag_{\qeightdd f}=\{(z,\qeightdd(f+q_\Sigma)(z)):z\in \qeightLdual\}$ of the differential of a
smooth function $f+q_\Sigma$ over the fixed Lagrangian plane $\qeightLdual$. It rounds
the corner while continuing to agree with $K$ along the four faces away from $v$;
it does not cap off $\Lambda_v$ by a compact surface, and it has no boundary
component on any contact sphere. There is therefore no compact filling to which a
Maslov or DGA obstruction could attach. The dichotomy ``$\Lambda_v$ bounds an
exact Lagrangian filling, or it is obstructed'' is not the dichotomy relevant to
local graphical smoothing; we are not in the filling category at all.

\paragraph{The graphical smoothing has vanishing Maslov class, unconditionally.}
Even granting the most generous reading of the obstruction --- as a constraint on
the Maslov class of the smoothed piece itself --- it is vacuous, because every
graph of a differential over a Lagrangian plane has trivial Maslov class. We
record this as a lemma.

\begin{qeightlemma}\label{lem:maslov-zero}
Let $\qeightLdual\subset\qeightR^4$ be a Lagrangian plane and $g\colon \qeightLdual\to\qeightR$ a
$C^2$ function. Then the graph
$\qeightLag=\{(z,\qeightdd g(z)):z\in \qeightLdual\}\subset \qeightLdual\oplus L\cong\qeightR^4$ has
vanishing Maslov class: the Maslov number of every loop in $\qeightLag$ is $0$.
\end{qeightlemma}

\begin{proof}
The Lagrangian Grassmannian $\Lambda(2)=U(2)/O(2)$ carries the universal Maslov
class, detected by the homotopy class of the map
$\det{}^2\colon\Lambda(2)\to S^1$ \cite[\S2.3]{McDuffSalamon}; the Maslov number
of a loop $\gamma$ in a Lagrangian submanifold is the degree of $\det{}^2$
composed with the Gauss map $z\mapsto T_z\qeightLag\in\Lambda(2)$ along $\gamma$.

Choose unitary coordinates so that $L=i\,\qeightLdual$ and $\qeightLdual=\qeightR^2\subset\qeightC^2$.
A plane that is the graph of a symmetric linear map $A\colon\qeightLdual\to L$ (i.e.\
transverse to $L$) is $\{x+iAx:x\in\qeightR^2\}$, represented in $U(2)/O(2)$ by the
unitary $(I+iA)(I+A^2)^{-1/2}$, whence
\[
\det{}^2\bigl(\mathrm{graph}\,A\bigr)
=\frac{\det(I+iA)}{\det(I-iA)}
=\prod_{k=1}^{2}\frac{1+i\lambda_k}{1-i\lambda_k},
\qquad \arg \det{}^2=2\sum_{k=1}^{2}\arctan\lambda_k,
\]
where $\lambda_1,\lambda_2$ are the (real) eigenvalues of the symmetric matrix
$A$. The graphical locus
$\mathcal G:=\{\mathrm{graph}\,A:A=A^{\mathsf T}\}\subset\Lambda(2)$ is, via
$A\mapsto\mathrm{graph}\,A$, homeomorphic to the linear space of symmetric
$2\times2$ matrices, hence contractible. Its continuous image
$\det{}^2(\mathcal G)\subset S^1$ is therefore null-homotopic in $S^1$ (a
continuous image of a contractible space supports no nontrivial loop class).

Now the tangent plane $T_z\qeightLag$ is the graph of the symmetric map
$\qeightHess g(z)\colon \qeightLdual\to L$ (the Hessian of $g$, symmetric by equality of
mixed partials), so the Gauss map of $\qeightLag$ factors through the contractible
graphical locus $\mathcal G$. Hence $\det{}^2\circ(\text{Gauss map})\colon\qeightLag\to
S^1$ factors through a contractible space, its degree on every loop is $0$, and the
Maslov class of $\qeightLag$ vanishes.
\end{proof}

Applying Lemma~\ref{lem:maslov-zero} with $g=f+q_\Sigma$ (smooth for
$f\in\qeightSm(\Sigma)$) shows the smoothed surface $\qeightLag_{\qeightdd f}$ has vanishing Maslov
class on every smoothing chart. Thus, far from being obstructed, the smoothed
Lagrangian is exact (it is the graph of an exact $1$-form $\qeightdd(f+q_\Sigma)$) and
Maslov-trivial; the putative obstruction is identically satisfied.

\paragraph{Where the four-valence enters, and what it is not.} The role of the
$4$-faces hypothesis is to guarantee, via Lemma~\ref{lem:linear}, that the vertex
admits the product-of-two-corners normal form, hence the homeomorphism
$\Sigma\cong\qeightLdual$ of Corollary~\ref{cor:Ldual} and the single-valued
piecewise-quadratic generating function $q_\Sigma=S(Q_1)+S(Q_2)$ with
$S(Q)=\tfrac12 Q|Q|$ (in the symplectic base coordinates $Q_i=(p_i-q_i)/\sqrt2$,
conjugate to $P_i=(p_i+q_i)/\sqrt2$, in which the corner reads $P_i=|Q_i|$). This
is a statement of linear symplectic geometry, not a parity or rotation-number
condition. In particular it is not a Maslov constraint: the obstruction-theoretic
quantities ($\mu$, $tb$, $r$, DGA augmentations) play no role because, as just
shown, the construction never leaves the Maslov-trivial graphical category. The
reconciliation of the two candidate answers is therefore clean: the
YES-construction is correct, and the NO-candidate rests on importing a
filling-category obstruction that does not govern local graphical smoothings.

\section*{Role of the $4$-faces-per-vertex hypothesis}

The hypothesis is used exactly once, in Lemma~\ref{lem:linear}: four Lagrangian
quadrants meeting at a vertex, with edge tangents $(v_1,v_2,u_1,u_2)$, force
$\omega(v_i,u_i)\ne0$ and hence a standard symplectic basis, which is what makes
the symplectic pairing $\Sigma\to\qeightLdual$ a homeomorphism
(Corollary~\ref{cor:Ldual}) and supplies the smooth structure used throughout.
With a different number of faces the local model changes and this normal form is
unavailable; the mechanism is the linear symplectic geometry of the vertex, not a
Maslov-class parity condition.

\section*{Remarks}

\begin{enumerate}\setlength\itemsep{4pt}
\item \textbf{Why smoothing functions, not Legendrian fillings.} The smoothing is
produced directly as graphs over $K$ via the conormal fibration, so no
holomorphic-curve or contact-fillability input is needed; the only analytic
ingredients are Weinstein's neighbourhood theorem and a partition-of-unity
estimate.
\item \textbf{The rescaling subtlety.} $\qeightSm(\Sigma)$ is closed under adding
smooth functions but not under scalar multiplication, so small smoothing
functions must be built by hand (Lemma~\ref{lem:small}); this is the one place a
genuine estimate is required. The smoothing path in Section~6 is therefore the
explicit family $f^\varepsilon\in\qeightSm(K)$ of Lemma~\ref{lem:cont-family}, which
stays inside the admissible class $\qeightSm(K)$ for every $\varepsilon>0$ and whose
$C^1$-norm tends to $0$; we do \emph{not} scale a smoothing function, since that
would leave $\qeightSm(K)$ (the coset $g+C^\infty(K)$ does not contain $0$).
\item \textbf{Matched normals.} Property (iii) of Definition~\ref{def:fibration}
ensures the two smooth structures on $K$ — from Corollary~\ref{cor:Ldual} and from
the edge collars — agree, so the global graph construction is consistent across
edges.
\item \textbf{Generalization.} The argument is local-to-global and uses only
linear symplectic normal forms plus partition-of-unity gluing, so it adapts to
other ambient symplectic manifolds once the vertex normal form is established.
\end{enumerate}

\begin{thebibliography}{9}
\bibitem{Hirsch} M.~W.~Hirsch.
  \emph{Differential Topology}. Graduate Texts in Mathematics, Springer, 1976.
\bibitem{McDuffSalamon} D.~McDuff and D.~Salamon.
  \emph{Introduction to Symplectic Topology}, 3rd ed.,
  Oxford University Press, 2017.
\end{thebibliography}

\endgroup
